\documentclass[11pt]{article}

\usepackage[margin=1in]{geometry}
\usepackage{amsmath,amssymb}
\usepackage{booktabs}
\usepackage{longtable}
\usepackage{enumitem}
\usepackage{hyperref}
\usepackage{xcolor}
\usepackage{listings}
\usepackage{microtype}
\usepackage[most]{tcolorbox}
\usepackage{titlesec}
\usepackage{array}

% Color definitions
\definecolor{solverblue}{RGB}{0,70,127}
\definecolor{warnred}{RGB}{180,30,30}
\definecolor{notebg}{RGB}{240,247,255}
\definecolor{warningbg}{RGB}{255,240,240}

\hypersetup{
  colorlinks=true,
  linkcolor=solverblue,
  citecolor=solverblue,
  urlcolor=solverblue
}

% Section formatting
\titleformat{\section}{\large\bfseries\color{solverblue}}{\thesection.}{0.6em}{}
\titleformat{\subsection}{\normalsize\bfseries}{\thesubsection.}{0.5em}{}

% Custom callout box
\tcbuselibrary{skins,breakable}
\newtcolorbox{notebox}[1][]{%
  colback=notebg, colframe=solverblue, fonttitle=\bfseries,
  title=#1, breakable, left=4pt, right=4pt
}
\newtcolorbox{warnbox}[1][]{%
  colback=warningbg, colframe=warnred, fonttitle=\bfseries,
  title=#1, breakable, left=4pt, right=4pt
}

\lstdefinestyle{terminal}{
  basicstyle=\ttfamily\small,
  columns=fullflexible,
  breaklines=true,
  frame=single,
  backgroundcolor=\color{black!3}
}

\lstdefinestyle{profiletree}{
  basicstyle=\ttfamily\footnotesize,
  columns=fullflexible,
  breaklines=false,
  frame=none,
  backgroundcolor=\color{black!3},
  xleftmargin=2em
}

\title{\textbf{LowMachReact-Hex}\\[0.4em]\large Technical Solver Report\\[0.2em]\normalsize Version~0.5.0}
\author{LowMachReact-Hex Development Notes}
\date{May 2026}

\begin{document}

\maketitle

\begin{abstract}
LowMachReact-Hex is a Fortran 2008 MPI finite-volume solver for hexahedral
meshes.  Its current implementation is best understood as a constant-density
low-Mach or incompressible projection solver with optional passive species
transport, optional passive sensible-enthalpy transport, and Cantera-backed
thermodynamic and transport-property evaluation.  The code provides a
hydrodynamic and thermodynamic foundation for later reacting-flow and radiation
coupling, but the current production path intentionally does not include
variable-density low-Mach coupling, reactions, reaction heat release,
species-diffusion enthalpy correction, or external radiation physics.

This report summarizes the solver architecture, data model, governing
equations, numerical algorithms, thermodynamic conventions, MPI strategy,
boundary-condition handling, profiling model, setup procedure, and known
limitations.
\end{abstract}

\tableofcontents

\section{Scope and Current Solver Interpretation}

LowMachReact-Hex currently solves constant-density incompressible flow with
optional passive scalar physics.  The pressure projection, momentum equation,
conservative face fluxes, CFL estimate, and diagnostic mass integral use the
configured flow density
\begin{equation}
  \rho_f = \rho_{\mathrm{params}}.
\end{equation}
In the code this value is represented by \texttt{params\%rho} and stored in the
transport-property container as \texttt{transport\%rho}.

Cantera may compute a thermodynamic density,
\begin{equation}
  \rho_{\mathrm{thermo}} = \rho_{\mathrm{Cantera}}(T,Y,p_0),
\end{equation}
but this value is diagnostic and future-use only.  It is stored as
\texttt{energy\%rho\_thermo} and is not used in the projection or momentum
equations.  A future variable-density low-Mach formulation would need a
different continuity constraint and pressure-correction derivation before this
thermodynamic density could participate in the flow solve.

\section{Enabled and Missing Physics}

\subsection{Enabled Physics}

The current staged implementation supports:
\begin{itemize}[leftmargin=*]
  \item constant flow and projection density;
  \item fractional-step pressure projection for incompressible flow;
  \item optional passive species mass-fraction transport;
  \item optional passive sensible-enthalpy transport;
  \item Fourier heat conduction driven by \(\nabla T\);
  \item Cantera evaluation of \(h(T,Y,p_0)\), \(T(h,Y,p_0)\), \(c_p\),
        \(\lambda\), mixture viscosity, species diffusivity, and diagnostic
        \(\rho_{\mathrm{thermo}}\);
  \item volumetric source storage through \texttt{qrad}, with
        \(q_{\mathrm{rad}}>0\) adding energy to the gas;
  \item VTU/PVTU output and CSV diagnostics;
  \item MPI-aware flat and nested profiling.
\end{itemize}

\subsection{Missing or Reserved Physics}

The following are not implemented in the current solver path:
\begin{itemize}[leftmargin=*]
  \item variable-density low-Mach divergence constraint;
  \item chemical reactions;
  \item reaction heat release;
  \item species-diffusion enthalpy correction,
        \[
          -\nabla \cdot \left(\sum_k h_k \mathbf{J}_k\right);
        \]
  \item external radiation physics;
  \item radiation spectral or wavenumber decomposition;
  \item active use of \(\rho_{\mathrm{thermo}}\) in the projection.
\end{itemize}

\section{Repository and Module Architecture}

\subsection{Source Layout}

The main solver source lives in \texttt{src/}.  The executable program is
\texttt{main.f90}, which coordinates input, mesh import, MPI setup, field
initialization, the timestep loop, diagnostics, output, and finalization.

\begin{longtable}{p{0.28\linewidth}p{0.64\linewidth}}
\toprule
\textbf{Module} & \textbf{Responsibility} \\
\midrule
\texttt{mod\_kinds} & Precision, constants, and common fatal-error handling. \\
\texttt{mod\_input} & Namelist parsing, parameter storage, and validation. \\
\texttt{mod\_mesh\_types} & Cell, face, patch, and mesh data structures. \\
\texttt{mod\_mesh\_io} & Native mesh-file reading. \\
\texttt{mod\_bc} & Boundary-condition parsing and face-state evaluation. \\
\texttt{mod\_fields} & Flow field allocation: velocity, pressure, divergence, face fluxes. \\
\texttt{mod\_flow\_projection} & Momentum prediction, pressure Poisson solve, velocity/flux correction. \\
\texttt{mod\_species} & Passive species advection, diffusion, correction, bounding, normalization. \\
\texttt{mod\_energy} & Sensible-enthalpy transport, temperature recovery, Cantera thermo sync. \\
\texttt{mod\_transport\_properties} & Constant or Cantera transport-property updates. \\
\texttt{mod\_mpi\_flow} & Replicated-mesh MPI ownership, exchange, gather, and reduction utilities. \\
\texttt{mod\_mpi\_radiation} & Reserved communicator/task structure for future radiation work. \\
\texttt{mod\_output} & VTU/PVTU and CSV diagnostics output. \\
\texttt{mod\_profiler} & Inclusive flat and nested MPI-aware timing reports. \\
\bottomrule
\end{longtable}

\subsection{Mesh Pipeline}

The mesh workflow separates geometry generation from the solver:
\begin{enumerate}[leftmargin=*]
  \item Gmsh creates hexahedral meshes with named physical boundary surfaces.
  \item A conversion tool reads the Gmsh mesh through \texttt{meshio}.
  \item The converter validates the current axis-aligned hexahedral assumptions.
  \item Native text files are written:
        \texttt{points.dat}, \texttt{cells.dat}, \texttt{faces.dat},
        \texttt{patches.dat}, and optionally \texttt{periodic.dat}.
  \item \texttt{mod\_mesh\_io} imports the native files into \texttt{mesh\_t}.
\end{enumerate}

\subsection{MPI Decomposition}

The current flow solver uses a replicated mesh and owned-cell decomposition.
Every MPI rank stores the full mesh, while \texttt{mod\_mpi\_flow} assigns each
rank a contiguous owned range of cell IDs.  Computation over cell-centered
unknowns is performed on owned cells, then synchronized by exchange or global
collective operations.

This design is intentionally conservative for the current problem sizes and
keeps the future radiation path simple: a spectral or ray-based radiation
decomposition can still access the full geometry on every rank.

\section{Primary Fields and State Variables}

\subsection{Flow Fields}

The flow state includes:
\begin{itemize}[leftmargin=*]
  \item cell-centered velocity \(\mathbf{u}\);
  \item previous velocity \(\mathbf{u}^n\);
  \item predicted velocity \(\mathbf{u}^{*}\);
  \item pressure \(p\);
  \item pressure correction \(\phi\);
  \item velocity divergence diagnostics;
  \item face-centered volumetric fluxes \(F_f\);
  \item previous explicit momentum RHS for Adams--Bashforth time stepping.
\end{itemize}

The stored face flux is oriented from owner cell to neighbor cell.  During
cell-local transport updates, it is reoriented so that \(F_{f,c}>0\) means
outflow from the currently updated cell.

\subsection{Species Fields}

Species are stored as mass fractions \(Y_k\) with array layout
\texttt{(species, cell)}.  The number of active species is determined from
\texttt{case.nml} or, in reserved reaction-discovery paths, from the Cantera
mechanism.  The present species equation is passive and contains no chemical
source term.

\subsection{Energy Fields}

The transported thermal state is mixture sensible enthalpy \(h\), not
temperature.  Temperature is a recovered thermodynamic variable:
\begin{equation}
  T = T(h,Y,p_0).
\end{equation}
The energy container also stores previous \(h\), previous \(T\), volumetric
source \(q_{\mathrm{rad}}\), \(c_p\), \(\lambda\), and diagnostic
\(\rho_{\mathrm{thermo}}\).

\section{Governing Equations}

\subsection{Incompressible Momentum and Continuity}

The incompressibility constraint is
\begin{equation}
  \nabla \cdot \mathbf{u} = 0.
\end{equation}
The momentum equation is advanced in fractional-step form:
\begin{equation}
  \frac{\partial \mathbf{u}}{\partial t}
  +
  \nabla \cdot (\mathbf{u}\mathbf{u})
  =
  -\frac{1}{\rho_f}\nabla p
  +
  \nabla \cdot (\nu \nabla \mathbf{u})
  +
  \mathbf{f}.
\end{equation}
Here \(\rho_f\) is the constant projection density, \(\nu\) is the kinematic
viscosity used by the flow solver, and \(\mathbf{f}\) is the configured body
acceleration.

\subsection{Passive Species}

For species \(k=1,\ldots,N_s\), the current model advances
\begin{equation}
  \frac{\partial Y_k}{\partial t}
  +
  \nabla \cdot (\mathbf{u}Y_k)
  =
  \nabla \cdot (D_k \nabla Y_k),
\end{equation}
with no reaction source:
\begin{equation}
  \dot{\omega}_k = 0.
\end{equation}

\subsection{Passive Sensible Enthalpy}

The current energy equation is
\begin{equation}
  \frac{\partial h}{\partial t}
  +
  \nabla \cdot (\mathbf{u}h)
  =
  \frac{1}{\rho_f}
  \nabla \cdot (\lambda \nabla T)
  +
  \frac{q_{\mathrm{rad}}}{\rho_f}.
\end{equation}
Conduction is driven by \(\nabla T\), not \(\nabla h\).  A positive
\(q_{\mathrm{rad}}\) increases gas sensible enthalpy.

\section{Finite-Volume Discretization}

For a hexahedral control volume \(c\), let \(V_c\) be the cell volume,
\(f \in \partial c\) a face, \(A_f\) the face area, \(d_f\) the normal distance
used for gradients, and \(\mathbf{n}_{f,c}\) the outward normal relative to
cell \(c\).

For a transported scalar \(\psi\), the solver uses upwind advection:
\begin{equation}
  \psi_f^{\mathrm{up}} =
  \begin{cases}
    \psi_c, & F_{f,c} \ge 0, \\
    \psi_{\mathrm{other}}, & F_{f,c} < 0,
  \end{cases}
\end{equation}
where \(\psi_{\mathrm{other}}\) is either the neighboring cell value or the
boundary state.

Central differences are used for diffusive face gradients in species and
thermal transport.  For temperature diffusion this has the form
\begin{equation}
  \left.\nabla T \cdot \mathbf{n}\right|_f
  \approx
  \frac{T_{\mathrm{other}} - T_c}{d_f}.
\end{equation}

\section{Projection Algorithm}

The projection method proceeds in three conceptual steps.

\subsection{Predictor}

An intermediate velocity \(\mathbf{u}^{*}\) is predicted without the new
pressure correction:
\begin{equation}
  \frac{\mathbf{u}^{*}-\mathbf{u}^{n}}{\Delta t}
  =
  \mathcal{R}^{n,n-1}
  -
  \frac{1}{\rho_f}\nabla p^n,
\end{equation}
where \(\mathcal{R}^{n,n-1}\) represents explicit advection, diffusion, and body
force terms.  The solver uses Adams--Bashforth 2 when previous RHS data are
available and falls back to Forward Euler on the first step.

\subsection{Pressure Correction}

The pressure correction potential \(\phi = p^{n+1}-p^n\) is obtained from the
constant-density Poisson equation
\begin{equation}
  \nabla^2 \phi =
  \frac{\rho_f}{\Delta t}\nabla \cdot \mathbf{u}^{*}.
\end{equation}
The matrix is applied in matrix-free form with cached geometric coefficients.
The linear system is solved by preconditioned conjugate gradient with a diagonal
Jacobi preconditioner.  For singular pure-Neumann systems, the implementation
pins the null space as needed.

\subsection{Correction}

The corrected velocity and pressure are
\begin{align}
  \mathbf{u}^{n+1}
  &=
  \mathbf{u}^{*}
  -
  \frac{\Delta t}{\rho_f}\nabla \phi, \\
  p^{n+1}
  &=
  p^n + \phi.
\end{align}
Face fluxes are corrected consistently with the pressure potential gradient so
that the updated velocity field is approximately divergence free.

\section{Species Transport}

Species transport uses explicit Euler time integration.  Each species receives
an advective flux and a diffusive flux.  The diffusive fluxes are corrected at
each face so that their sum is zero:
\begin{equation}
  \sum_k J_{k,f}^{\mathrm{corr}} = 0.
\end{equation}
After the explicit update, mass fractions are bounded to \([0,1]\) and
renormalized locally so that
\begin{equation}
  \sum_k Y_k = 1.
\end{equation}
This keeps the passive species vector usable for Cantera composition calls and
prevents small numerical drift from accumulating in the mixture sum.

\section{Energy and Cantera Thermodynamics}

\subsection{Sensible-Enthalpy Convention}

When Cantera thermodynamics are enabled, the stored enthalpy is not Cantera's
absolute formation-including enthalpy.  The solver stores sensible enthalpy
relative to a reference temperature:
\begin{equation}
  h(T,Y,p_0)
  =
  h_{\mathrm{abs}}^{\mathrm{Cantera}}(T,Y,p_0)
  -
  h_{\mathrm{abs}}^{\mathrm{Cantera}}(T_{\mathrm{ref}},Y,p_0).
\end{equation}
This convention avoids artificial heat release in non-reacting passive mixing.

\subsection{Temperature Recovery}

To recover temperature from stored sensible enthalpy, the code reconstructs an
absolute enthalpy target:
\begin{equation}
  h_{\mathrm{target,abs}}
  =
  h
  +
  h_{\mathrm{abs}}^{\mathrm{Cantera}}(T_{\mathrm{ref}},Y,p_0),
\end{equation}
then uses Cantera's HPY inversion to find \(T\).

\subsection{Option A Coupling Rule}

The code follows the Option A convention:
\begin{quote}
  \(h\) is the transported thermodynamic state.
\end{quote}
When species transport changes \(Y\), the solver preserves \(h\) and recovers
the temperature consistent with the new composition:
\begin{equation}
  T^{*} = T(h,Y^{n+1},p_0).
\end{equation}
The solver must not recompute
\begin{equation}
  h^{*} = h(T^n,Y^{n+1},p_0)
\end{equation}
from the old temperature and new composition, because that would add or remove
sensible enthalpy without a corresponding energy flux or source.

\subsection{Thermo Sync}

The preferred Cantera energy-path synchronization is a combined operation:
\begin{equation}
  (T,c_p,\lambda,\rho_{\mathrm{thermo}})
  =
  \operatorname{sync}(h,Y,p_0).
\end{equation}
The sync must preserve the transported \(h\).  If a Cantera call writes a
temporary enthalpy value as part of property evaluation, the Fortran path saves
and restores the transported field.

The parameter \texttt{thermo\_update\_interval} is reserved and must remain
equal to 1 in supported runs.  Internal caching can reduce redundant Cantera
work, but the logical thermodynamic state must remain synchronized every energy
step.

\section{Runtime Ordering}

The main timestep ordering is:
\begin{enumerate}[leftmargin=*]
  \item Update CFL diagnostics or dynamic timestep state when needed.
  \item Refresh transport properties on \texttt{transport\_update\_interval}.
        Cantera may update \(\mu\) and \(D_k\); flow density remains
        \(\rho_f=\texttt{params\%rho}\).
  \item Advance momentum and pressure projection.
  \item Advance passive species, if enabled.
  \item Advance passive sensible enthalpy, if enabled.
  \item Recover or synchronize temperature and thermodynamic properties.
  \item Write diagnostics and VTU/PVTU output on output steps.
\end{enumerate}

For species-enabled Cantera thermo runs, a pre-flux energy sync is needed after
species transport because \(Y\) may have changed.  For species-disabled Cantera
thermo runs, that pre-flux sync may be skipped because composition is unchanged
and the previous post-flux sync remains valid.  A post-flux sync is performed
after \(h\) is updated.

\section{Boundary Conditions}

Boundary patches are named in the mesh and configured through \texttt{case.nml}.
The code supports a legacy patch type plus field-specific overrides:
\begin{itemize}[leftmargin=*]
  \item velocity boundary type;
  \item pressure boundary type;
  \item species boundary type;
  \item temperature/enthalpy boundary type.
\end{itemize}

For fixed-temperature boundaries in Cantera thermo mode, boundary enthalpy is
computed from the boundary temperature and boundary composition:
\begin{equation}
  h_b = h(T_b,Y_b,p_0).
\end{equation}
When species are enabled, \(Y_b\) comes from the species boundary condition.
This is important for hot or cold fuel and oxidizer inlets, where defaulting to
the interior or bath-gas composition would make boundary enthalpy inconsistent.

\section{Transport Properties and Cantera Bridge}

\texttt{mod\_transport\_properties} supports constant fallback properties and
Cantera-backed transport properties.  Cantera may compute:
\begin{itemize}[leftmargin=*]
  \item dynamic viscosity \(\mu\);
  \item species diffusivities \(D_k\);
  \item thermal conductivity \(\lambda\) in the energy thermo path;
  \item \(c_p\) in the energy thermo path;
  \item diagnostic thermodynamic density.
\end{itemize}

When variable viscosity is disabled, the flow solver keeps the configured
constant viscosity for validation even if Cantera species diffusivities are
enabled.  When variable viscosity is enabled, the kinematic viscosity is
computed from Cantera \(\mu\) and the constant flow density:
\begin{equation}
  \nu = \frac{\mu}{\rho_f}.
\end{equation}

The C++ bridge maps solver species names to Cantera mechanism species names,
constructs full Cantera mass-fraction vectors, fills missing bath gas when
needed, and caches repeated thermo or transport evaluations conservatively.

\section{Diagnostics, Output, and Profiling}

The solver writes global diagnostics and ParaView-compatible visualization
files.  Output may include velocity, pressure, divergence, temperature,
enthalpy, \(q_{\mathrm{rad}}\), \(c_p\), \(\lambda\), thermal diffusivity,
\(\rho_{\mathrm{thermo}}\), species mass fractions, species diffusivities, and
cell IDs.

The profiler reports inclusive wall-clock timing.  Flat profiler rows are not
additive when nested timers are enabled.  Common top-level timers include:
\begin{itemize}[leftmargin=*]
  \item \texttt{Transport\_Update};
  \item \texttt{Projection\_Step};
  \item \texttt{Species\_Transport};
  \item \texttt{Energy\_Transport};
  \item \texttt{Diagnostics\_Write\_Flow};
  \item \texttt{Diagnostics\_Write\_Energy};
  \item \texttt{Output\_Write\_VTU}.
\end{itemize}
Cantera energy sync timers include \texttt{Energy\_Cantera\_PreSync} and
\texttt{Energy\_Cantera\_PostSync}.

\section{Build and Run Procedure}

The project expects a Conda environment created from \texttt{environment.yml}
and an available OpenMPI compiler/runtime.

\begin{lstlisting}[style=terminal]
conda env create -f environment.yml
conda activate react_env

# Load OpenMPI if your system uses environment modules.
module load openmpi

# Build.
make BUILD=release

# Run.
mpirun -np 4 ./lowmach_react_hex cases/channel_flow/case.nml
\end{lstlisting}

The Makefile creates a compiled binary under \texttt{build/} and a top-level
\texttt{lowmach\_react\_hex} launcher.  Users should run the launcher.  It
prepends the active Conda library directory so Cantera and the C++ runtime are
resolved from the environment.

\section{Validation Strategy}

Recommended validation proceeds from low-risk to high-risk configurations:
\begin{enumerate}[leftmargin=*]
  \item Build in debug mode and run a zero- or few-step startup smoke test.
  \item Validate single-rank incompressible flow on a simple case.
  \item Repeat the same case with multiple MPI ranks and compare diagnostics.
  \item Enable passive species transport with constant diffusivity.
  \item Enable constant-property energy transport.
  \item Enable Cantera thermo without transported species.
  \item Enable species-coupled Cantera thermo and verify the Option A behavior.
  \item Compare profiler output before and after performance changes.
\end{enumerate}

For each validation run, keep the case file, rank count, build mode, terminal
output, diagnostics CSV files, and representative visualization output.

\section{Known Limitations and Future Work}

The most important current limitation is the deliberate constant-density
projection formulation.  Future variable-density low-Mach support should not be
added by simply replacing \(\rho_f\) with \(\rho_{\mathrm{thermo}}\); it requires
a redesigned divergence constraint and pressure equation.

Additional future work includes:
\begin{itemize}[leftmargin=*]
  \item reaction source terms and reaction heat release;
  \item species-diffusion enthalpy correction;
  \item radiation physics coupled through \(q_{\mathrm{rad}}\);
  \item improved mesh partitioning beyond replicated full-mesh storage;
  \item higher-order bounded advection for scalar transport;
  \item broader regression tests across energy/species/Cantera modes.
\end{itemize}


\section{Performance Characterization}

The solver's performance is dominated by the coupling between thermodynamic
property synchronization and the iterative pressure correction solver. Based
on a 16{,}000-step benchmark on a multi-rank hexahedral mesh (total wall time
83.6~s avg across ranks), the following granular profiling tree was collected.

\subsection{Full Nested Profiling Tree}

\begin{lstlisting}[style=profiletree]
Region Tree                          Calls    Avg(s)    Avg%
------------------------------------------------------------
Total_Simulation                         1   83.564   100.00
  MPI_Communication                      1    0.005     0.01
  Transport_Update                      16    0.186     0.22
    Transport_Setup                     16    0.000     0.00
    Transport_Cantera_Call              16    0.074     0.09
    Transport_Exchange_Diff             16    0.112     0.13
      MPI_Communication                 16    0.112     0.13
  Projection_Step                    16000   33.994    40.68
    Projection_Momentum_RHS          16000    1.732     2.07
    Projection_AB2                   16000    0.169     0.20
    Projection_Predict_Flux          16000    2.105     2.52
      PredictFlux_Compute            16000    0.408     0.49
      PredictFlux_Balance            16000    1.194     1.43
      MPI_Communication              16000    0.771     0.92
    Projection_Poisson_RHS           16000    0.133     0.16
    Projection_PCG                   16000   28.053    33.57
      MPI_Communication            6629753   17.609    21.07
      Pressure_Matvec              2199251    5.786     6.92
    Projection_Pressure_Update       16000    0.074     0.09
    Projection_Correction            16000    1.370     1.64
  Species_Transport                  16000    1.531     1.83
    MPI_Communication                16000    0.209     0.25
  Energy_Transport                   16000   47.180    56.46
    Energy_Cantera_PreSync           16000   21.847    26.14
    Energy_PreFlux_Exchange          16000    3.794     4.54
      MPI_Communication              32000    3.773     4.52
    Energy_Flux_Update               16000    0.653     0.78
    Energy_Cantera_PostSync          16000   18.591    22.25
    Energy_Final_Exchange            16000    2.029     2.43
      MPI_Communication              48000    2.002     2.40
\end{lstlisting}

\subsection{Sub-Module Cost Summary}

\begin{longtable}{p{0.32\linewidth}p{0.1\linewidth}p{0.1\linewidth}p{0.38\linewidth}}
\toprule
\textbf{Sub-Region} & \textbf{Avg~s} & \textbf{\%} & \textbf{Technical Observation} \\
\midrule
\texttt{Energy: Cantera PreSync} & 21.85 & 26.1\% & HPY inversion per cell, pre-flux. \\
\texttt{Energy: Cantera PostSync} & 18.59 & 22.3\% & Property refresh after $h$ update. \\
\texttt{Projection: PCG MPI} & 17.61 & 21.1\% & Global dot products and halo exchanges. \\
\texttt{Projection: Matvec} & 5.79 & 6.9\% & Matrix-free Laplacian (efficient). \\
\texttt{Energy: PreFlux Comm} & 3.77 & 4.5\% & MPI exchange of $T$, $h$, $c_p$, $\lambda$. \\
\texttt{Energy: Final Comm} & 2.00 & 2.4\% & Post-update thermodynamic field sync. \\
\texttt{Projection: Predictor} & 2.11 & 2.5\% & Momentum RHS and flux computation. \\
\texttt{Species: Transport+Comm} & 1.53 & 1.8\% & Explicit advection-diffusion (low cost). \\
\bottomrule
\end{longtable}

\subsection{Deep-Dive Analysis}

\subsubsection{Thermodynamic Bottleneck (48.4\% total)}

The \texttt{Energy\_Cantera\_Sync} cycle is the single dominant cost. Pre-flux
state recovery (HPY inversion) consumes \textbf{26.1\%} and post-flux property
refreshment consumes \textbf{22.3\%} of total wall time. The cost is nearly
symmetric because both operations call the same per-cell Cantera HPY solve.
This thermodynamic compute cost is approximately \textbf{20$\times$} the cost
of the entire hydrodynamic predictor step.

\begin{notebox}[Optimization Opportunity: Cantera]
Batch cell-state HPY inversions across the mesh rather than calling Cantera
cell-by-cell. Alternatively, amortize the cost by caching thermo properties
when $T$ and $Y$ changes are below a tolerance threshold.
\end{notebox}

\subsubsection{PCG Communication Sensitivity (21.1\% total)}

The Pressure Poisson solver (PCG) is executed 16{,}000 times with a total
of 6{,}629{,}753 \texttt{MPI\_Communication} calls. Of the \textbf{33.6\%}
total time in \texttt{Projection\_PCG}, over \textbf{63\%} (21.1\% of total)
is pure MPI overhead. This identifies the solver as \emph{interconnect-latency
bound}: the bottleneck is the global dot products required by the CG
algorithm, not the arithmetic.

\begin{notebox}[Optimization Opportunity: PCG]
Replace the standard PCG with a pipelined or CA-CG variant to overlap
communication with computation. Alternatively, coarsen the problem with
AMG to reduce iteration count and thus global reduction frequency.
\end{notebox}

\subsubsection{Transport Exchange Disparity}

Species transport incurs only \textbf{0.25\%} MPI overhead, whereas energy
transport incurs \textbf{6.9\%} (PreFlux + Final exchanges combined). The
cause is the number of fields that must be synchronized: species exchanges
only $Y_k$, while energy exchanges require $T$, $h$, $c_p$, and $\lambda$
for boundary-cell consistency with Cantera.

\subsection{Recommended Optimization Priority}

\begin{enumerate}[leftmargin=*]
  \item \textbf{Batch Cantera HPY Inversions}: Profile-directed batching of
        per-cell thermo calls is expected to yield the largest wall-time
        reduction (target: reduce 48.4\% cost by $\sim$50\%).
  \item \textbf{Pipelined or Communication-Avoiding PCG}: Reduces the
        frequency of global reductions, which currently account for 63\%
        of the Poisson solve time.
  \item \textbf{Fused Energy Exchange}: Combine the PreFlux and Final MPI
        exchanges into a single non-blocking operation where possible.
\end{enumerate}

\section{Summary}

LowMachReact-Hex is currently a constant-density projection solver with
well-defined passive species and sensible-enthalpy transport extensions.
The critical state convention is:
\begin{itemize}[leftmargin=*]
  \item $h$ (sensible enthalpy relative to $T_{\mathrm{ref}}$) is the
        transported thermodynamic scalar.
  \item $T$ is a \emph{recovered} variable, obtained via Cantera HPY
        inversion from $(h,Y,p_0)$.
  \item $\rho_{\mathrm{thermo}}$ (Cantera thermodynamic density) is
        \emph{diagnostic only} and is never used in the projection.
\end{itemize}

These conventions make the solver internally consistent while preserving
a clear path toward future extensions. The performance analysis identifies
Cantera thermodynamic synchronization (48.4\%) and PCG MPI communication
(21.1\%) as the two primary targets for optimization before scaling to
larger meshes or longer runs.

\begin{warnbox}[Future Architecture Warning]
Do not activate variable-density low-Mach coupling by simply substituting
$\rho_{\mathrm{thermo}}$ for $\rho_f$ in the projection. This requires a
full redesign of the divergence constraint and pressure-correction equation.
\end{warnbox}

\end{document}
